% generator_I05.tex -- Prop I.5 (Theor.): isosceles triangle -- base angles equal.
% Geometry and text from jemmybutton (CC-BY-SA). Build: lualatex generator_I05.tex
\documentclass[12pt]{article}
\usepackage{geometry}\geometry{a5paper, margin=1.5cm}
\usepackage[lining]{ebgaramond}
\usepackage{amssymb}
\usepackage{byrne}
\def\mpPre{u := 1cm; textLabels := false;}

\begin{document}

\defineNewPicture{%
  pair A, B, C, D, E;%
  A := (0, 0);%
  B := A shifted (u, -2u);%
  C := B xscaled -1;%
  D := 9/5[A,B];%
  E := 9/5[A,C];%
  byAngleDefine(B, A, C, byblack, SOLID_SECTOR);%
  byAngleDefine(A, B, C, byblue,   SOLID_SECTOR);%
  byAngleDefine(B, C, A, byblue,   SOLID_SECTOR);%
  byAngleDefine(C, B, E, byyellow, SOLID_SECTOR);%
  byAngleDefine(D, C, B, byyellow, SOLID_SECTOR);%
  byAngleDefine(B, D, C, byred,    SOLID_SECTOR);%
  byAngleDefine(C, E, B, byred,    SOLID_SECTOR);%
  byAngleDefine(E, B, D, byblack,  ARC_SECTOR);%
  byAngleDefine(D, C, E, byblack,  ARC_SECTOR);%
  draw byNamedAngleResized(BAC,ABC,BCA,CBE,DCB,BDC,CEB);%
  byLineDefine(B, D, byyellow, SOLID_LINE, REGULAR_WIDTH);%
  byLineDefine(C, E, byyellow, SOLID_LINE, REGULAR_WIDTH);%
  byLineDefine(B, E, byblue,   SOLID_LINE, REGULAR_WIDTH);%
  byLineDefine(C, D, byblue,   SOLID_LINE, REGULAR_WIDTH);%
  byLineDefine(A, B, byred,    SOLID_LINE, REGULAR_WIDTH);%
  byLineDefine(A, C, byred,    SOLID_LINE, REGULAR_WIDTH);%
  byLineDefine(B, C, byblack,  SOLID_LINE, REGULAR_WIDTH);%
  draw byNamedLineSeq(0)(CD,noLine,BC,noLine,BE,CE,AC,AB,BD);%
  draw byLabelsOnPolygon(E, C, A, B, D, C, B)(ALL_LABELS, 0);%
}

% STATEMENT
\noindent
\begin{minipage}[t]{0.45\linewidth}
\centering\vspace{0pt}%
\drawCurrentPicture
\end{minipage}\hfill
\begin{minipage}[t]{0.52\linewidth}
\vspace{0pt}%
\textbf{Book~I.\enspace Prop.\ V. Theor.}

\medskip\noindent\itshape
In any isosceles triangle \drawLine[bottom]{BC,AC,AB}
if the equal sides be produced, the external angles at the
base are equal, and the internal angles at the base are
also equal.

\bigskip

% PROOF
\noindent\upshape
Produce \drawUnitLine{AB} and \drawUnitLine{AC}~(post.~II),\\
take $\drawUnitLine{BD} = \drawUnitLine{CE}$~(pr.~I.3);\\
draw \drawUnitLine{BE} and \drawUnitLine{CD}.

\medskip
Then in
\offsetPicture{8pt}{0pt}{\drawFromCurrentPicture{%
  startAutoLabeling;%
  draw byNamedAngle(BAC);%
  draw byNamedLineSeq(0)(BE,CE,AC,AB);%
  stopAutoLabeling;%
}} and
\offsetPicture{8pt}{0pt}{\drawFromCurrentPicture{%
  startAutoLabeling;%
  draw byNamedAngle(BAC);%
  draw byNamedLineSeq(0)(BD,CD,AC,AB);%
  stopAutoLabeling;%
}}\\
we have $\drawUnitLine{AB,BD} = \drawUnitLine{AC,CE}$~(const.),\\
\drawAngle{BAC} common to both,\\
and $\drawUnitLine{AB} = \drawUnitLine{AC}$~(hyp.)\\
$\therefore \drawAngle{BCA,DCB} = \drawAngle{ABC,CBE}$,
$\drawUnitLine{BE} = \drawUnitLine{CD}$
and $\drawAngle{CEB} = \drawAngle{BDC}$~(pr.~I.4).

\medskip
Again in \drawLine{BE,CE,BC} and \drawLine{BD,CD,BC}\\
we have $\drawUnitLine{BD} = \drawUnitLine{CE}$,
$\drawAngle{CEB} = \drawAngle{BDC}$
and $\drawUnitLine{BE} = \drawUnitLine{CD}$,\\
$\therefore \drawAngle{DCE,DCB} = \drawAngle{EBD,CBE}$
and $\drawAngle{DCB} = \drawAngle{CBE}$~(pr.~I.4)\\
but $\drawAngle{BCA,DCB} = \drawAngle{ABC,CBE}$,
$\therefore \drawAngle{BCA} = \drawAngle{ABC}$~(ax.~III).

\begin{flushright}Q.\ E.\ D.\end{flushright}
\end{minipage}

\end{document}
